\section{The string-flow framework}\label{sec:framework}

\subsection{Words, weights, and molecules}

\begin{definition}[Valid word]\label{def:wordValid}
A word $w=(t_0,\dots,t_{L-1})$ is \emph{valid from} $x_0$ if the
recursive definition of $\wordOrbit(w,x_0)$ is exact, i.e.\
\begin{equation*}
  (5x_i+1)\equiv 0 \pmod{2^{t_i}}
\end{equation*}
for every prefix state $x_i$.
\end{definition}

\begin{definition}[Prefix weight]\label{def:wordWeight}
For a word $w=(t_0,\dots,t_{L-1})$, the prefix weight after $j$ steps
is
\begin{equation*}
  \wordWeight_j(w)=\sum_{i<j} t_i,
\end{equation*}
and the total weight is $\wordWeight(w)=\wordWeight_L(w)$.
\end{definition}

\begin{definition}[Word molecule]\label{def:wordA}
For a word $w$ of length $L$, the \emph{word molecule} is the unique
integer $\wordA(w)$ satisfying
\begin{equation*}
  2^{\wordWeight(w)}\,\wordOrbit(w,x_0)
    = 5^L\,x_0 + \wordA(w)
\end{equation*}
for every initial state $x_0$.  Equivalently,
\begin{equation*}
  \wordA(w)=\sum_{j=0}^{L-1}
    2^{\wordWeight_j}\,5^{L-1-j}.
\end{equation*}
Equivalently, \(\wordA([])=0\) and
\(\wordA(w{::}t)=5\,\wordA(w)+2^{\wordWeight(w)}\);
this recursion is the definition used in Lean.
\end{definition}

The word molecule converts orbit arithmetic into exact prefix
arithmetic.  It is the object that appears in the Lean statements
\texttt{StringFlow.Word.wordA} and \texttt{StringFlow.PMI.aTotal5}.
The formalisation is written in Lean~\cite{lean} using
Mathlib~\cite{mathlib}.

\begin{theorem}[Exact word orbit identity]\label{thm:word-identity}
For every valid word $w=(t_0,\dots,t_{L-1})$ from $x_0$,
\begin{equation*}
  2^{\wordWeight(w)}\,\wordOrbit(w,x_0)
    =5^L\,x_0+\wordA(w).
\end{equation*}
\end{theorem}

\begin{proof}[Proof sketch]
Proceed by induction on the length of $w$.  For the singleton word
$[t]$,
\begin{equation*}
  2^t\,\wordOrbit([t],x_0)=5x_0+1,
\end{equation*}
and $\wordA([t])=1$.  For $w'=w{::}t$,
\begin{equation*}
  2^{\wordWeight(w')}\,\wordOrbit(w',x_0)
    =2^{\wordWeight(w)}\,(5\,\wordOrbit(w,x_0)+1)
    =5\,(2^{\wordWeight(w)}\,\wordOrbit(w,x_0))+2^{\wordWeight(w)},
\end{equation*}
which matches the recursive formula for $\wordA(w')$.  The identity is
the exact ledger used throughout the paper.
\end{proof}

\subsection{The valuation decomposition}

Every positive integer $n$ has a unique decomposition
\begin{equation*}
  n=2^{\vtwo{n}}\,\operatorname{oddpart}(n),
\qquad \operatorname{oddpart}(n)\equiv1\pmod2.
\end{equation*}
In Lean this is recorded by \texttt{n\_eq\_two\_pow\_mul\_oddPart} and
\texttt{oddPart\_odd\_of\_pos}.  The exact step weight
$t=\vtwo{5x+1}$ is the largest weight for which $(5x+1)/2^t$ is an
integer.  Word validity only asks for divisibility, so a legal word may
use any $t\le\vtwo{5x+1}$.  This is the precise sense in which
$\GeneralOrbit$ is weaker than $\FullOrbit$.

\subsection{Valuation toolkit}

The following elementary facts are used throughout without comment:
\begin{itemize}[leftmargin=*]
\item $n=2^{\vtwo{n}}\operatorname{oddpart}(n)$ with
  $\operatorname{oddpart}(n)$ odd;
\item if $u$ is odd, then $\vtwo{2^k u}=k$;
\item for odd $x$, $\vtwo{5x+1}\ge1$;
\item $T(x)\le(5x+1)/2$;
\item a C3 step from $x\ge1$ satisfies $(5x+1)/2^t<x$.
\end{itemize}
These elementary facts are formalised in Lean: the odd-part
decomposition, the valuation rules, the step bound, and the C3
decrease lemma.

\begin{proposition}[Exact orbit bound]\label{prop:orbit-bound}
For every $n\ge2$,
\begin{equation*}
  \Orbit_7(n)<5^n.
\end{equation*}
This is proved in Lean
(\texttt{CycleBridge.fiveXPlusOneOrbit\_lt\_five\_pow}).
\end{proposition}

\begin{proof}
For $n=2$ the claim is checked directly.  If
$x=\Orbit_7(k)<5^k$, then
\begin{equation*}
  5x+1<5^{k+1}+1<2\cdot5^{k+1},
\end{equation*}
so $(5x+1)/2<5^{k+1}$.  The accelerated step satisfies
$T(x)\le(5x+1)/2$, hence $\Orbit_7(k+1)<5^{k+1}$.
\end{proof}

\subsection{The first steps}

The first 16 states and step weights make the definitions concrete:
\begin{center}
\scriptsize
\begin{tabular}{@{}ccccccccccccccccc@{}}
\toprule
$n$ & 0 & 1 & 2 & 3 & 4 & 5 & 6 & 7 & 8 & 9 & 10 & 11 & 12 & 13 & 14 & 15 \\
\midrule
$\Orbit_7(n)$ & 7 & 9 & 23 & 29 & 73 & 183 & 229 & 573 & 1433 & 3583 & 4479 & 5599 & 6999 & 8749 & 21873 & 54683 \\
$t_n$ & 2 & 1 & 2 & 1 & 1 & 2 & 1 & 1 & 1 & 2 & 2 & 2 & 2 & 1 & 1 & 3 \\
\bottomrule
\end{tabular}
\end{center}
The step from depth 15 to depth 16 has weight $3$, and the full orbit
then reaches $34177$.  The table is proved by explicit rewriting in
\texttt{FinitePrefix.lean}.

\begin{example}\label{ex:wordA}
The word $w=[2,1]$ is valid from $7$:
\begin{equation*}
  7 \xrightarrow{t=2} 9 \xrightarrow{t=1} 23,
\end{equation*}
so $\wordOrbit(w,7)=23$.  Its prefix weights are
$\wordWeight_0=0$, $\wordWeight_1=2$, $\wordWeight=3$, and
\begin{equation*}
  \wordA(w)=2^0\cdot5^1+2^2\cdot5^0=9.
\end{equation*}
The exact word-orbit identity gives
\begin{equation*}
  2^3\cdot23=184=5^2\cdot7+9.
\end{equation*}
The word $[1,2]$ is valid from $9$:
\begin{equation*}
  9 \xrightarrow{t=1} 23 \xrightarrow{t=2} 29,
\qquad
  \wordA([1,2])=2^0\cdot5^1+2^1\cdot5^0=7,
\end{equation*}
and $2^3\cdot29=232=5^2\cdot9+7$.

Validity is weaker than tracking the accelerated orbit: the word
$[1,1]$ is valid from $9$ in the divisibility sense, because
$46/2=23$ and $116/2=58$ are integers, but the exact second weight at
$23$ is $v_2(116)=2$, not $1$.  Hence
$\wordOrbit([1,1],9)=58$ is an intermediate state not on the
accelerated orbit of $7$.  This is precisely the distinction between
$\GeneralOrbit$ and $\FullOrbit$ used in \Cref{def:reach}.
\end{example}

\subsection{The prefix ledger}

Let $\wordWeight_n$ be the prefix weight after $n$ exact steps from
$7$, and let $\wordA_n$ be the word molecule of that exact prefix word.
The first values are:
\begin{center}
\scriptsize
\begin{tabular}{@{}rrrr@{}}
\toprule
$n$ & $t_n$ & $\wordWeight_n$ & $\wordA_n$ \\
\midrule
0 & -- & 0 & 0 \\
1 & 2 & 2 & 1 \\
2 & 1 & 3 & 9 \\
3 & 2 & 5 & 53 \\
4 & 1 & 6 & 297 \\
5 & 1 & 7 & 1549 \\
6 & 2 & 9 & 7873 \\
7 & 1 & 10 & 39877 \\
8 & 1 & 11 & 200409 \\
9 & 1 & 12 & 1004093 \\
10 & 2 & 14 & 5024561 \\
11 & 2 & 16 & 25139189 \\
12 & 2 & 18 & 125761481 \\
13 & 2 & 20 & 629069549 \\
14 & 1 & 21 & 3146396321 \\
15 & 1 & 22 & 15734078757 \\
16 & 3 & 25 & 78674588089 \\
\bottomrule
\end{tabular}
\end{center}

Each row is an instance of the exact word-orbit identity:
\begin{equation*}
  2^{\wordWeight_n}\,\Orbit_7(n)=5^n\cdot7+\wordA_n.
\end{equation*}
For example, $2^{25}\cdot34177=5^{16}\cdot7+78674588089$.  The column
$\wordA_n$ is the word molecule of the exact prefix word and is the
value of \texttt{StringFlow.Word.wordA} on that prefix.

\subsection{The PMI identity}

\begin{theorem}[PMI prefix identity]\label{thm:pmi}
Let $w$ be a valid word of length $L$ from $x_0$, let
$\wordWeight_j$ be its prefix weights, and define the margin
\begin{equation*}
  m_j=j\log_2 5-\wordWeight_j.
\end{equation*}
Then
\begin{equation*}
  \sum_{j=0}^{L-1} 2^{-m_j}
    =\frac{5\,\wordA(w)}{5^L}.
\end{equation*}
\end{theorem}

The cleared-denominator form
\begin{equation*}
  \sum_{j=0}^{L-1} 2^{\wordWeight_j}\,5^{L-j}
    =5\,\wordA(w)
\end{equation*}
is the identity \(\mathrm{aTotal5}=5\,\mathrm{aTotal}\) formalised in
\texttt{Pmi.lean}.  For a closed cycle word of period $P$ and cycle
state $m$, substituting the cycle equation gives
\begin{equation*}
  \mathrm{aTotal5}(P,t)=5m(2^T-5^P),
\qquad T=\wordWeight(w).
\end{equation*}

\begin{proof}[Derivation sketch]
Write
\begin{equation*}
  \wordA(w)=\sum_{j=0}^{L-1}
    2^{\wordWeight_j}\,5^{L-1-j}.
\end{equation*}
Then
\begin{equation*}
  \frac{5\,\wordA(w)}{5^L}
    =\sum_{j=0}^{L-1}2^{\wordWeight_j}\,5^{-j}.
\end{equation*}
Since $2^{-m_j}=2^{\wordWeight_j}5^{-j}$, the identity follows.  The
calculation is a finite rational identity; no limits or estimates are
involved.
\end{proof}

\subsection{Prefix budgets and PMI-B}

The PMI-B layer counts the positions in a word that violate the prefix
budget.  For a closed cycle word of period $P$, the total weight
satisfies
\begin{equation*}
  2^{\wordWeight(w)}\,m=5^P\,m+\wordA(w)
\end{equation*}
for the cycle state $m$, so $\wordWeight(w)>P\log_2 5$.  The PMI-B
inequalities then compare the number of ``bad'' prefixes with the
available budget.  The corresponding PMI-B no-bad-prefix theorem and
the margin-balance theorem are compiled in Lean.

\begin{lemma}[PMI-B bad-prefix count]\label{lem:pmi-b-count}
For a closed cycle word of period $P$, the number of bad prefixes and
the total weight satisfy the PMI-B bounds
\begin{equation*}
  \operatorname{bad}(w)\le
  \frac{2^{\wordWeight(w)}}{5^P}\,\wordA(w),
\end{equation*}
and every true prefix of a small-budget word satisfies the margin
balance inequality.  The counting bound and the no-bad-prefix
conclusion are compiled in Lean.
\end{lemma}

\begin{lemma}[Margin balance]\label{lem:margin-balance}
Let $\alpha=\log_2 5$, and let $C_1,C_2,C_3$ be the number of steps of
weight $1$, weight $2$, and weight at least $3$ in a closed cycle word.
Then the nonnegative-margin prefix condition implies
\begin{equation*}
  (3-\alpha)C_3\le(\alpha-1)C_1+(\alpha-2)C_2,
\end{equation*}
and the strict version replaces $\le$ by $<$.  The margin-balance
inequalities are compiled in Lean.
\end{lemma}

\subsection{PMI-B derivation}

Write $B=\operatorname{bad}(w)$ for the number of bad proper prefixes
of a closed cycle word of period $P$, and let $T$ be its total weight.
The $j=0$ term of $\mathrm{aTotal5}$ contributes $5^P$, and every bad
prefix $j$ contributes at least another $5^P$:
\begin{equation*}
  5^P+B\cdot5^P\le\mathrm{aTotal5}(P,t).
\end{equation*}
Together with the cycle equation
\begin{equation*}
  \mathrm{aTotal5}(P,t)=5m(2^T-5^P),
\end{equation*}
this is the exact counting bound
\begin{equation*}
  5^P+B\cdot5^P\le5m(2^T-5^P).
\end{equation*}

\begin{corollary}[Small budget, no bad prefix]\label{cor:pmi-b-no-bad}
If $5m(2^T-5^P)<2\cdot5^P$, then
\begin{equation*}
  2^{\wordWeight_j}<5^j\qquad(1\le j<P).
\end{equation*}
\end{corollary}

\begin{proof}
The counting bound gives $B\cdot5^P<5^P$, hence $B=0$.
\end{proof}

The statements are \texttt{pmi\_b\_count} and
\texttt{pmi\_b\_no\_bad\_prefix} in \texttt{Pmi.lean}.

\subsection{PMI-B worked example}

The word $[3,1]$ has prefix weights
\begin{equation*}
  \wordWeight_0=0,\quad \wordWeight_1=3,\quad \wordWeight_2=4,
\end{equation*}
and molecule $\wordA([3,1])=13$.  Its margins are
\begin{equation*}
  m_0=0,\qquad
  m_1=\log_2 5-3<0,\qquad
  m_2=2\log_2 5-4>0,
\end{equation*}
so the proper prefix of length $1$ is bad and $B=1$.  The cleared PMI
sum is
\begin{equation*}
  \mathrm{aTotal5}(2,t)
    =5^2+2^3\cdot5^1=25+40=65,
\end{equation*}
and the counting bound reads
\begin{equation*}
  5^2+B\cdot5^2=25+25=50<65.
\end{equation*}
The exact orbit word $[2,1,2]$ from $7$, by contrast, has all margins
positive.  This is the exact arithmetic that the PMI-B layer counts:
one bad prefix costs one extra $5^P$, and no estimate replaces the
count.

\subsection{Rise steps and residues modulo 5}

A rise step of weight $1$ satisfies $2y=5x+1$, hence $y\equiv3
\pmod5$.  A rise step of weight $2$ satisfies $4y=5x+1$, hence
$y\equiv4\pmod5$.  These two residue rules are the only constraints
that a rise step imposes on the next state modulo $5$; they are
formalised in \texttt{rise\_step\_next\_mod\_five} and are used by the
cycle-bridge block-head residue lemma.

\subsection{Block equations}

A block is a contiguous segment of the orbit.  For a block starting at
state $r_0$ with step weights $u_1,\dots,u_L$, the block equation is
\begin{equation*}
  2^{U}\,r_L=5^L\,r_0+B,
\end{equation*}
where $U=\sum u_i$ and $B$ is the block molecule.  The reset block of
weight $t\in\{1,2\}$ and parameter $\delta\in\{1,3\}$ satisfies
\begin{equation*}
  2^t(r_j+1)=5^{k_0+1}s+\delta5^j
\end{equation*}
in the $t=2$ case, and the corresponding $t=1$ equation
\begin{equation*}
  2(r_j+1)=5^{k_0+1}s+5^j-2.
\end{equation*}
These are the equations encoded in
\texttt{S6Audit.ResetHeadEq} and used by the corrected window bound.

\begin{example}[Block equation]\label{ex:block}
The block starting at $29$ with weights $[1,1,2]$ has
$L=3$, $U=4$, and block molecule $B=39$:
\begin{equation*}
  2^4\cdot229=3664=5^3\cdot29+39.
\end{equation*}
\end{example}

\subsection{Lean correspondence}

\begin{center}
\begin{tabular}{lll}
\toprule
Mathematical object & Lean name & Status \\
\midrule
Word validity & \texttt{Word.wordValid} & compiled \\
Word orbit & \texttt{Word.wordOrbit} & compiled \\
Word molecule & \texttt{Word.wordA} & compiled \\
Prefix weight & \texttt{wordWeight} & compiled \\
PMI identity & \texttt{Pmi.aTotal5} & compiled \\
PMI-B budget & \texttt{cycleWord\_pmi\_b\_count} & compiled \\
Exact orbit bound & \texttt{fiveXPlusOneOrbit\_lt\_five\_pow} & compiled \\
Full orbit reachability & \texttt{FullOrbitFrom7} & compiled \\
General orbit reachability & \texttt{GeneralOrbitFrom7} & compiled \\
Reset window reachability & \texttt{ResetWindowReachability} & compiled \\
\bottomrule
\end{tabular}
\end{center}

\subsection{Reachability: full orbit versus general orbit}

\begin{definition}[Full and general orbit reachability]\label{def:reach}
Let $\FullOrbit(x)$ mean that $x$ lies on the exact accelerated orbit
of $7$:
\begin{equation*}
  \FullOrbit(x)\iff \exists n,\ \Orbit_7(n)=x.
\end{equation*}
Let $\GeneralOrbit(x)$ mean that $x$ is reachable from $7$ by some
finite word whose divisions are exact, without any restriction on the
weights.
\end{definition}

The difference is essential.  The general orbit permits submaximal
weights and therefore intermediate states that are not states of the
exact accelerated orbit; the full orbit does not.  The 25-state table
used in earlier regression work is a finite check and is not used as a
substitute for the full-orbit exclusion.

\subsection{Regression versus proof}

\begin{center}
\small
\begin{tabular}{@{}p{0.42\textwidth}p{0.52\textwidth}@{}}
\toprule
Regression data & Proof requirement \\
\midrule
25-state table & 17-state exact expansion \\
simulated orbit values & explicit rewriting, no \texttt{native\_decide} \\
statistical weight profile & exact word ledger with valuation identities \\
heuristic unboundedness & contradiction from a cycle assumption \\
\bottomrule
\end{tabular}
\end{center}

Regression tables are used in this paper only to locate candidates;
they never enter a proof step.  The 25-state table, for example, can
confirm that no early state repeats, but it cannot rule out a later
cycle.  Every proof step instead uses the exact word-orbit identity,
the explicit finite prefix, or a compiled Lean statement.  This
distinction is the reason the paper separates $\FullOrbit$ from
$\GeneralOrbit$ and keeps the 17-state expansion explicit.

\begin{definition}[Reset window reachability]\label{def:reset-window}
The predicate $\ResetWindow(j,k_0,t,\delta,s)$ asserts the existence of
$r$ and $r_j$ such that
\begin{equation*}
  k_0+1\le j,\qquad s5^{k_0}=r+1,\qquad
  \GeneralOrbit(r),\qquad \FullOrbit(r_j),
\end{equation*}
and the exact reset equation $\mathrm{ResetHeadEq}(s,j,k_0,t,\delta,r_j)$
holds, together with the size bound $s<5^{j-1-k_0}$, the oddness of
$s$ and $r_j$, and $5\nmid s$.
\end{definition}

The predicate is the reachability input of the corrected window bound.
It fixes the same $j,k_0,t$ as the window under consideration; a weaker
predicate that only fixes the difference $j-k_0-1$ would allow the
reset witness to use different indices.

\subsection{Notation discipline}

The step weight at depth $j$ is $t_j=\vtwo{5\Orbit_7(j)+1}$.  The
prefix weight before that step is $\wordWeight_{j-1}$.  These two
quantities are not interchangeable: $t_j$ is the cost of the next step,
while $\wordWeight_{j-1}$ is the accumulated cost before it.  The proof
never substitutes one for the other.

\subsection{Glossary}

\begin{center}
\small
\begin{tabular}{@{}p{0.26\textwidth}p{0.68\textwidth}@{}}
\toprule
Term & Meaning \\
\midrule
word & finite list of step weights \\
valid word & word whose orbit divisions are exact \\
prefix weight & accumulated weight before a step \\
word molecule & exact prefix sum of a word \\
margin & $j\log_2 5-\wordWeight_j$ \\
block & contiguous orbit segment \\
rise step & step of weight $1$ or $2$ \\
C3 step & step of weight at least $3$ \\
reset window & reachable block with exact reset equation \\
failure window & reset block whose valuation contradicts the bound \\
full orbit & exact accelerated orbit of $7$ \\
general orbit & orbit with arbitrary legal weights \\
\bottomrule
\end{tabular}
\end{center}

\subsection{Variable dictionary}

\begin{center}
\small
\begin{tabular}{@{}p{0.14\textwidth}p{0.80\textwidth}@{}}
\toprule
Symbol & Meaning \\
\midrule
$j$ & reset step depth \\
$k_0$ & exponent in $s5^{k_0}=r+1$ \\
$t$ & reset step weight, $t\in\{1,2\}$ \\
$\delta$ & reset parameter: $1$ for $t=1$, $\{1,3\}$ for $t=2$ \\
$s$ & odd terminal part with $5\nmid s$ \\
$e$ & incoming step weight, $e\ge2$ \\
$g$ & full-orbit state $\Orbit_7(j-1)$ \\
$x$ & candidate predecessor $2^{e-1}g+\delta5^{j-1}$ \\
$r_j$ & block head after the reset step \\
$r_s$ & terminal state of the block tail \\
$t_j$ & step weight at depth $j$ \\
$\wordWeight_j$ & prefix weight at depth $j$ \\
$L$ & block-tail length \\
$H_s$ & unified-core threshold parameter \\
\bottomrule
\end{tabular}
\end{center}

\subsection{Why loose estimates do not close the problem}

Heuristic models predict roughly logarithmic growth, but the orbit
contains long stretches of weight-$2$ steps that consume 2-adic value,
interspersed with weight-$1$ and weight-$\ge3$ steps.  Bounds that only
control the average growth are too weak: the decisive obstruction is an
exact inequality with constants $13$ and $+12$, and any weakening of
these constants destroys the contradiction.  Probabilistic equidistribution
statements and transcendence estimates do not provide the exact
2-adic cancellations used in the unified core.

A heuristic word with the same average weight can fail the exact
valuation identity by a single unit of 2-adic cancellation.  The
arithmetic counterexample in the divergence bridge is exactly such a
witness: it satisfies every size, parity, and coprimality condition
but is not orbit-reachable.  The proof therefore keeps every
divisibility condition exact.
